\documentclass[11pt]{article}
\usepackage{fltcourse}

\title{Level 7: The classification of mod $3$ hardly ramified representations}
\author{Kevin Buzzard}
\date{}

\begin{document}
\maketitle

\noindent
In the previous level we reduced Fermat's Last Theorem to the conjunction of
three statements $B_{6a}$ (lifting), $B_{6b}$ (spreading out) and $B_{6c}$
(classification at $3$). In this level we \emph{prove} $B_{6c}$ outright, and so
lop it off the list: the boss statement is
\[
  B_7 := B_{6a}\wedge B_{6b},
\]
and the boss theorem says that $B_7$ implies FLT. Everything here is known in
the 1980s or earlier; the level is long only because the pieces must be
assembled with care.

\medskip\noindent
Throughout, $\chi$ denotes the $3$-adic cyclotomic character and $\chibar$ its
reduction modulo $3$. We freely use the notions from Level 5 (coefficient
ring, hardly ramified, flat at a prime, decomposition/inertia/Frobenius). The
heart of the matter is the following classification of mod $3$ hardly ramified
representations, whose proof occupies most of the level.

\begin{lemma}[Mod $3$ classification]\label{lem:mod3}
Let $k$ be a finite field of characteristic $3$, let $V$ be a two-dimensional
$k$-vector space, and let $\rhobar\colon\GQ\to\Aut_k(V)$ be a hardly ramified
representation. Then, for a suitable basis,
\[
  \rhobar\;\cong\;\begin{pmatrix}\chibar & * \\ 0 & 1\end{pmatrix}.
\]
In particular, there are no irreducible mod $3$ hardly ramified
representations.
\end{lemma}

We first show how Lemma~\ref{lem:mod3} yields $B_{6c}$ (and hence the boss),
and then prove the lemma.

\section{Two stability properties}

The reduction of $B_{6c}$ to Lemma~\ref{lem:mod3} needs two facts about how
the hardly ramified condition behaves. Both were implicit in Level 5; we make
them explicit.

\begin{proposition}[Base change of coefficients]\label{prop:basechange}
Let $R_1\to R_2$ be a continuous homomorphism of coefficient rings of the same
odd residue characteristic $\ell$. If $\rho_1\colon\GQ\to\Aut_{R_1}(\calV_1)$ is
hardly ramified, then so is its base change
$\rho_2:=\rho_1\otimes_{R_1}R_2$ on $\calV_2:=\calV_1\otimes_{R_1}R_2$.
\end{proposition}

\begin{proof}
The determinant of $\rho_2$ is the image of $\det\rho_1=\chi$ under
$R_1\to R_2$; by uniqueness of the map $\Zl\to R_2$ this is again $\chi$, giving
(HR1). The kernel of $\rho_2$ contains that of $\rho_1$, so $\rho_2$ is
unramified wherever $\rho_1$ is, giving (HR2). For (HR4), tensoring the
$\Gal(\Qbar_2/\Q_2)$-stable line $\calW_1$ and its quotient (all free) with
$R_2$ preserves the required shape, the unramified character $\psi$ being
unchanged. Finally (HR3): given an open ideal $J_2\subseteq R_2$, its preimage
$J_1\subseteq R_1$ is open, and $R_1/J_1\hookrightarrow R_2/J_2$ is an injection
of finite rings; the finite Galois module $\calV_2/J_2\calV_2$ is obtained from
the flat module $\calV_1/J_1\calV_1$ by writing $R_2/J_2$ as a quotient of a
finite free $R_1/J_1$-module and using that flatness is preserved by finite
sums and quotients. Hence $\rho_2$ is flat at $\ell$.
\end{proof}

\begin{proposition}[Change of lattice]\label{prop:lattice}
Let $\calO$ be the ring of integers of a finite extension $K/\Ql$ ($\ell$
odd), let $\rho\colon\GQ\to\Aut_\calO(\calV)$ be hardly ramified, and put
$V:=\calV\otimes_\calO K$. If $\calV'\subseteq V$ is any $\GQ$-stable
$\calO$-lattice \textup{(}free of rank $2$\textup{)}, then the associated
representation $\rho'\colon\GQ\to\Aut_\calO(\calV')$ is again hardly ramified.
\end{proposition}

\begin{proof}
Conditions (HR1) and (HR2) depend only on $V=\calV\otimes K=\calV'\otimes K$, so
hold for $\rho'$. For (HR3), scaling $\calV'$ by a power of $\ell$ we may assume
$\ell^N\calV\subseteq\calV'\subseteq\calV$; then for any open ideal $J$ the
module $\calV'/J\calV'$ is a subquotient of $\calV/\ell^N J\calV$, which is
flat as $\ell^N J$ is open, so $\calV'/J\calV'$ is flat. For (HR4), let
$\calW\subseteq\calV$ be the given $\Gal(\Qbar_2/\Q_2)$-stable line; then
$W:=\calW\otimes K\subseteq V$ is $\Gal(\Qbar_2/\Q_2)$-stable with $V/W$ acted on
through $\psi$, and $\calW':=W\cap\calV'$ is a $\Gal(\Qbar_2/\Q_2)$-stable line
in $\calV'$ whose quotient $\calV'/\calW'$ injects into $V/W\cong K$, hence is
free of rank one with $\Gal(\Qbar_2/\Q_2)$ acting through $\psi$.
\end{proof}

\section{From the mod $3$ lemma to $B_{6c}$ and the boss}

\begin{proposition}\label{prop:B6c}
Lemma~\ref{lem:mod3} implies Statement $B_{6c}$.
\end{proposition}

\begin{proof}
Let $\rho\colon\GQ\to\Aut_\calO(\calV)$ be hardly ramified, with $\calO$ the
integers of a finite extension $K/\Q_3$; write $\m$ for its maximal ideal and
$k=\calO/\m$ (a finite field of characteristic $3$), and put
$V=\calV\otimes_\calO K$ and $\tilde\rho=\rho\otimes_\calO K$.

By Proposition~\ref{prop:basechange} the reduction $\rhobar=\rho\bmod\m$ is
hardly ramified over $k$, so by Lemma~\ref{lem:mod3} it is reducible, with
$\rhobar^{\mathrm{ss}}\cong\chibar\oplus\mathbf1$ (and $\chibar\neq\mathbf1$).

\emph{$\tilde\rho$ is reducible.} Suppose not. By Ribet's Lemma (see
J.-P.~Serre's report, or Ribet's 1976 paper \emph{A modular construction of
unramified $p$-extensions of $\Q(\mu_p)$}), an irreducible $\tilde\rho$ with
$\rhobar^{\mathrm{ss}}$ a sum of two \emph{distinct} characters admits, for each
of the two orderings, a $\GQ$-stable lattice whose reduction is a
\emph{nonsplit} extension in that order. Taking the order opposite to the one
provided by Lemma~\ref{lem:mod3}, we obtain a stable lattice $\calV'$ with
$\calV'/\m\calV'\cong\left(\begin{smallmatrix}\mathbf1&*\\0&\chibar\end{smallmatrix}\right)$,
nonsplit. By Proposition~\ref{prop:lattice} the lattice representation $\rho'$
is hardly ramified, hence so is its reduction $\calV'/\m\calV'$; but a hardly
ramified $k$-representation of this shape contradicts Lemma~\ref{lem:mod3}
(which forces $\chibar$ on the stable line). So $\tilde\rho$ is reducible.

\emph{Conclusion.} As $\tilde\rho$ is reducible, some $K$-basis of $V$ puts it
in the form $\left(\begin{smallmatrix}\alpha_1&*\\0&\alpha_2\end{smallmatrix}\right)$
with characters $\alpha_1,\alpha_2\colon\GQ\to K^\times$; choosing a stable
lattice inside $V$ we may take $\rho$ itself upper-triangular over $\calO$. The
restrictions $\alpha_i|_{\Gal(\Qbar_3/\Q_3)}$ are subquotients of the flat
representation $\rho|_{\Gal(\Qbar_3/\Q_3)}$, hence flat. Because $3$ is odd, the
ramification index $e=1$ of $\Q_3/\Q_3$ satisfies $e<3-1$, and in this range
Fontaine's classification of $1$-dimensional flat representations (from his
1977 Ast\'erisque memoir; cf.\ Raynaud, 1974) applies: \emph{the only flat
$1$-dimensional representations of $\Gal(\Qbar_3/\Q_3)$ are the unramified
characters and the unramified twists of the cyclotomic character.} The
condition (HR4) at $2$ forces both eigenvalues of $\tilde\rho$ on an inertia
group $I_2$ to be $1$, so $\alpha_1,\alpha_2$ are unramified at $2$; being also
unramified away from $2,3$, they are unramified outside $3$. By the
Kronecker--Weber theorem an abelian character unramified outside $3$ factors
through $\Gal(\Q(\mu_{3^\infty})/\Q)$, on which inertia at $3$ acts through all
of $\Z_3^\times$; there is thus no room for an unramified twist, and each
$\alpha_i$ is either trivial or equal to $\chi$. Their product is
$\det\rho=\chi$, so exactly one of them is $\chi$. By Lemma~\ref{lem:mod3}
applied to the reduction, the trivial character is the quotient and $\chi$ the
sub, giving $\rho\cong\left(\begin{smallmatrix}\chi&*\\0&1\end{smallmatrix}\right)$.
This is $B_{6c}$.
\end{proof}

\begin{theorem}[Boss of Level 7]\label{thm:boss7}
Statements $B_{6a}$ and $B_{6b}$ together imply Fermat's Last Theorem, modulo
results known in the 1980s.
\end{theorem}

\begin{proof}
By Lemma~\ref{lem:mod3} and Proposition~\ref{prop:B6c}, Statement $B_{6c}$
holds unconditionally (modulo the 1980s). Hence $B_7=B_{6a}\wedge B_{6b}$
implies $B_{6a}\wedge B_{6b}\wedge B_{6c}=B_6$, which by the boss of the
previous level implies FLT.
\end{proof}

It remains to prove Lemma~\ref{lem:mod3}.

\section{Discriminants and ramification}

We shall bound the number field cut out by a mod $3$ hardly ramified
representation, using its discriminant. Recall that a number field $L$ of
degree $n$ has a \emph{discriminant} $\disc(L)\in\Z$, the determinant of the
trace form $(e_i,e_j)\mapsto\tr_{L/\Q}(e_ie_j)$ on a $\Z$-basis of the ring of
integers. We record the classical facts we need, all pre-1980.

\begin{itemize}
\item A prime $p$ divides $\disc(L)$ iff it ramifies in $L$, i.e.\ iff the
image of an inertia group $I_p$ in $\Gal(L/\Q)$ is nontrivial.
\item \emph{(Tame discriminant formula.)} If $L/\Q$ is Galois and $p$ is
\emph{tamely} ramified with ramification index $e$ (the order of the image of
$I_p$), then the exact power of $p$ dividing $\disc(L)$ is
$p^{\,[L:\Q](e-1)/e}$. \textup{(}This follows from
$v_{P}(\mathcal D_{L/\Q})=e-1$ for the different, and $\disc=N(\mathcal
D)$.\textup{)}
\item \emph{(Structure of inertia.)} An inertia group $I_p$ has a normal
pro-$p$ subgroup $P_p$ \textup{(}wild inertia\textup{)} with pro-cyclic
quotient $I_p/P_p\cong\prod_{q\neq p}\Z_q$, and any Frobenius acts on $I_p/P_p$
by $x\mapsto x^p$.
\item \emph{(Fontaine's bound.)} If $k$ has characteristic $p$ and
$\rhobar\colon\Gal(L/\Q)\hookrightarrow\Aut_k(V)$ is flat at $p$, then the
exact power $p^{b}$ of $p$ dividing $\disc(L)$ satisfies
$b<\dfrac{p}{p-1}[L:\Q]$. \textup{(}Th\'eor\`eme~A of Fontaine, \emph{Il n'y a
pas de vari\'et\'e ab\'elienne sur $\Z$}, 1985.\textup{)}
\item \emph{(Odlyzko--Poitou discriminant bounds.)} There are explicit lower
bounds for the \emph{root discriminant} $\disc(L)^{1/[L:\Q]}$ of a number field
in terms of its degree. We use two consequences \textup{(}from Poitou,
\emph{Sur les petits discriminants}, 1977, and Odlyzko\textup{)}: a totally
imaginary field with $[L:\Q]\geq6$ has root discriminant $>4.557$; and a field
with root discriminant $<8.25$ has $[L:\Q]\leq 14$.
\end{itemize}

\section{Proof of the mod $3$ classification (Lemma~\ref{lem:mod3})}

Let $k$ be finite of characteristic $3$ and $\rhobar\colon\GQ\to\Aut_k(V)$
hardly ramified, $V$ of rank $2$. Since $V$ is finite with the discrete
topology, $\rhobar$ factors through an injection
$\rhobar\colon\Gal(L/\Q)\hookrightarrow\Aut_k(V)$ for a finite Galois extension
$L/\Q$; write $G=\Gal(L/\Q)$. By (HR2), $L$ is unramified outside $2$ and $3$,
so $\disc(L)=\pm2^{a}3^{b}$ for some $a,b\geq0$. By (HR1) and the fact that
$\det\rhobar=\chibar$ is the mod $3$ cyclotomic character, $L$ contains
$\Q(\zeta_3)=\Q(\sqrt{-3})$; hence $[L:\Q]$ is even and $L$ is totally
imaginary.

\subsection*{Local analysis at $2$}
By (HR4), $\rhobar|_{\Gal(\Qbar_2/\Q_2)}\cong\left(\begin{smallmatrix}\psi\chibar&*\\0&\psi\end{smallmatrix}\right)$
with $\psi$ unramified and $\chibar$ unramified at $2$, so the image of an
inertia group $I_2$ lies in the unipotent subgroup $\left(\begin{smallmatrix}1&*\\0&1\end{smallmatrix}\right)\cong k^{+}$,
an elementary abelian $3$-group. As $I_2/P_2$ is pro-cyclic and $P_2$ (pro-$2$)
maps trivially to a $3$-group, the image of $I_2$ is cyclic of order dividing
$3$; thus $L$ is tame at $2$ with $e_2\in\{1,3\}$. By the tame discriminant
formula, $a\in\{0,\ 2[L:\Q]/3\}$.

\subsection*{Local analysis at $3$: the tame case}
Let $J_3\subseteq G$ be the image of an inertia group $I_3$. Suppose first that
$3\nmid|J_3|$ (tame case). Then $J_3$, a finite quotient of the pro-cyclic
$I_3/P_3$, is cyclic of order prime to $3$; over a finite extension $k'$ of $k$
it is diagonalisable, $\rhobar|_{J_3}\cong\phi_1\oplus\phi_2$ with
$\phi_i\colon J_3\to k'^\times$. A Frobenius $\Frob_3$ acts on $J_3$ by
$x\mapsto x^3$, and conjugating $\rhobar|_{J_3}$ by $\rhobar(\Frob_3)$ gives an
isomorphic representation; therefore
$\phi_1\oplus\phi_2\cong\phi_1^{3}\oplus\phi_2^{3}$. Hence
$\{\phi_1,\phi_2\}=\{\phi_1^3,\phi_2^3\}$, so in either case $\phi_i^9=\phi_i$,
i.e.\ $\phi_i^{8}=1$. Thus $|J_3|\mid 8$ (in fact $|J_3|\leq 8$), and by the
tame formula $b\leq 7[L:\Q]/8$.

Combining with $a\leq 2[L:\Q]/3$, the root discriminant satisfies
\[
  \disc(L)^{1/[L:\Q]}\;\leq\;2^{2/3}\,3^{7/8}\;<\;4.16.
\]
Since $L$ is totally imaginary, the Odlyzko--Poitou bound forces $[L:\Q]<6$;
being even, $[L:\Q]\in\{2,4\}$. Then $G$ is abelian of order prime to $3$, so
over some finite extension $k''$ of $k$ we have
$\rhobar\cong\chi_1\oplus\chi_2$ with characters $\chi_i$. Each $\chi_i$ is
unramified at $2$ (inertia at $2$ has order $1$ or $3$, but $|G|$ is prime to
$3$, so it is $1$) and tame at $3$ (its order divides $4$); being unramified
outside $3$ and tame at $3$, it factors through
$\Gal(\Q(\zeta_3)/\Q)$, so $\chi_i\in\{\mathbf1,\chibar\}$. Their product is
$\det\rhobar=\chibar$, so $\rhobar\cong\mathbf1\oplus\chibar$ over $k''$; by
Brauer--Nesbitt the same holds over $k$, giving
$\rhobar\cong\left(\begin{smallmatrix}\chibar&*\\0&1\end{smallmatrix}\right)$
with $*=0$. This settles the tame case.

\subsection*{Local analysis at $3$: the wild case}
Now suppose $3\mid|J_3|$ (wild case). Since $\rhobar$ is flat at $3$,
Fontaine's bound gives $b<\tfrac{3}{2}[L:\Q]$, whence
\[
  \disc(L)^{1/[L:\Q]}\;<\;2^{2/3}\,3^{3/2}\;<\;8.25 .
\]
By the Odlyzko--Poitou bound, $[L:\Q]\leq14$. But $[L:\Q]$ is even
($\Q(\zeta_3)\subseteq L$) and divisible by $3$ (as $3\mid|J_3|$ divides
$|G|=[L:\Q]$); hence $[L:\Q]\in\{6,12\}$.

Set $H:=\Gal(L/\Q(\zeta_3))$, an index-$2$ (hence normal) subgroup of $G$, of
order $3$ or $6$. A group of order $3$ or $6$ has a \emph{unique} subgroup $P$
of order $3$; being unique it is characteristic in $H$, hence normal in $G$.
Now $\rhobar$ is a faithful mod $3$ representation and $P$ has order $3$, so $P$
acts nontrivially and unipotently on $V$; a $3$-group acting on a two-dimensional
$k$-vector space with nonzero, proper fixed space fixes exactly a $k$-line. As
$P$ is normal in $G$, this line is $G$-stable, so
$\rhobar\cong\left(\begin{smallmatrix}\phi_1&*\\0&\phi_2\end{smallmatrix}\right)$
where $\phi_1$ is the character on the stable line. Both $\phi_i$ have order
prime to $3$ (their values lie in $k^\times$), are unramified at $2$ (by (HR4),
as above) and tame at $3$; hence, exactly as in the tame case, each is $\mathbf1$
or $\chibar$, with product $\chibar$. So one is $\chibar$ and the other
$\mathbf1$.

If $\phi_1=\chibar$ and $\phi_2=\mathbf1$ we are done. It remains to rule out
$\phi_1=\mathbf1$, $\phi_2=\chibar$, i.e.\
$\rhobar\cong\left(\begin{smallmatrix}\mathbf1&*\\0&\chibar\end{smallmatrix}\right)$.
In that case $H=\Gal(L/\Q(\zeta_3))$ acts through
$\left(\begin{smallmatrix}1&*\\0&1\end{smallmatrix}\right)$ (as $\chibar|_H=1$),
so its faithful image is a $3$-group; since $|H|\in\{3,6\}$ this forces $|H|=3$
and $|G|=6$. Then the image of a decomposition group $D_3$ is all of $G$, so the
image of $\Gal(\Qbar_3/\Q_3)$ in $\Aut_k(V)$ has order divisible by $3$, and
$\rhobar|_{\Gal(\Qbar_3/\Q_3)}$ is a nonsplit, ramified extension of $\chibar$
by $\mathbf1$ which is flat at $3$. A theorem of Ribet (in the 1976 paper cited
above), resting on Raynaud's uniqueness of finite flat models when $e<\ell-1$
(here $\ell=3$, $e=1$), shows that no such flat extension exists. This
contradiction rules out the case $\phi_1=\mathbf1$, and completes the wild
case.

\medskip
In both cases $\rhobar\cong\left(\begin{smallmatrix}\chibar&*\\0&1\end{smallmatrix}\right)$,
proving Lemma~\ref{lem:mod3}. \hfill$\square$
\end{document}
